\FloatBarrier
\section{پیش‌پردازش}
\label{sec:preproc}

سیگنال‌های خام را نمی‌توان مستقیم به مدل داد. \lr{EEG} پر از نویز برق شهر و حرکت چشم است، \lr{EMG} خام نوسانی سریع حول صفر است و نه شدت فعالیت، و حسگرهای نیرو گاهی نیرویی را ثبت می‌کنند که از دست شرکت‌کننده نیست. پیش‌پردازش این مشکل‌ها را رفع می‌کند و همه‌ی سیگنال‌ها را به شکلی درمی‌آورد که مدل بتواند از آن‌ها یاد بگیرد.

پیش‌پردازش دو مرحله دارد:
\begin{enumerate}
\item \textbf{مرحله‌ی مشترک} (شکل~\ref{fig:preproc-shared}): از فایل‌های خام تا سیگنال‌های تمیز ۵۰۰ هرتز. هر دو مدل از خروجی همین مرحله استفاده می‌کنند. نتیجه برای هر سری یک بار ساخته و در یک \lr{cache} ذخیره می‌شود.
\item \textbf{مرحله‌ی ویژه‌ی \lr{v2}} (شکل~\ref{fig:preproc-v2}): از \lr{cache} تا ورودی‌ها و هدف ۵۰ هرتز مدل \lr{v2}. این مرحله هر بار که \lr{v2} داده را بارگذاری می‌کند اجرا می‌شود و چیزی روی دیسک نمی‌نویسد.
\end{enumerate}

\subsection{داده‌ی خام و تقسیم سری‌ها}

برای هر شرکت‌کننده ۹ سری اصلی وجود دارد (فایل‌های \lr{HS\_P\{p\}\_S\{s\}.mat}). یک سری دیگر هم هست (\lr{ST}) که در هیچ‌کدام از دو مدل استفاده نمی‌شود. در هر فایل سه نوع سیگنال هست:
\begin{itemize}
\item \lr{EEG}: ۳۲ کانال با ۵۰۰ هرتز؛
\item \lr{EMG}: ۵ عضله با ۴۰۰۰ هرتز؛
\item سینماتیک: ۳۶ ستون با ۵۰۰ هرتز، شامل مکان چهار حسگر (جسم، نوک انگشت اشاره، نوک شست و مچ، به سانتی‌متر) و نیروها و گشتاورهای دو صفحه‌ی نیروسنج روی جسم.
\end{itemize}

سری‌ها سه نوع‌اند: ۵ سری که فقط وزن جسم در آن‌ها عوض می‌شود (\lr{W})، ۲ سری که فقط جنس سطح عوض می‌شود (\lr{S}) و ۲ سری که هر دو عوض می‌شوند (\lr{B}). سری‌های ۷ تا ۹ همیشه از نوع \lr{W} هستند. تقسیم برای هر شرکت‌کننده با یک قاعده‌ی ثابت انجام می‌شود:
\begin{itemize}
\item \textbf{آزمون (۲ سری):} دومین سری \lr{W} از میان سری‌های ۱ تا ۶ و دومین سری \lr{B}؛
\item \textbf{\lr{validation} (۱ سری):} سری ۹؛
\item \textbf{آموزش (۶ سری):} بقیه، یعنی اولین \lr{W} از ۱ تا ۶، سری‌های ۷ و ۸، هر دو سری \lr{S} و اولین \lr{B}.
\end{itemize}
با این قاعده، آزمون هم سری‌ای دارد که فقط وزن در آن عوض می‌شود و هم سری‌ای که وزن و سطح هر دو عوض می‌شوند، و آموزش از هر سه نوع نمونه دارد. یک قاعده‌ی مهم در همه‌ی مراحل: هر آماره‌ای که برای نرمال کردن لازم است (میانگین، انحراف معیار، صدک ۹۵) فقط از سری‌های آموزش حساب می‌شود. اگر از سری‌های آزمون هم استفاده می‌شد، اطلاعات آزمون به آموزش نشت می‌کرد.

\begin{figure}[p]
\centering
\resizebox{\textwidth}{!}{% Shared preprocessing: raw HS_P*_S*.mat -> 500 Hz cache (both models).
\begin{latin}
\begin{tikzpicture}[node distance=4mm, font=\footnotesize, box/.append style={text width=3.9cm}]
% EEG column
\node[io, text width=3.9cm] (e0) {EEG, 32 channels, 500 Hz};
\node[ebox, below=of e0] (e1) {band-pass 0.5--40 Hz\\(FIR, Hamming, zero-phase)};
\node[ebox, below=of e1] (e2) {notch 50 Hz (FIR)};
\node[ebox, below=of e2] (e3) {ICA (FastICA, per series)\\remove components that follow Fp1\\(eye-movement proxy, $|z|>2$)};
\node[ebox, below=of e3] (e4) {second band-pass 0.5--40 Hz};
\node[ebox, below=of e4] (e5) {robust clip per channel\\$\pm\max(5\,\mathrm{MAD}/0.6745,\ 50)$};
% EMG column
\node[io, right=6mm of e0, text width=3.9cm] (m0) {EMG, 5 muscles, 4000 Hz};
\node[box, below=of m0, fill=intc!8, draw=intc!70!black] (m1) {band-pass 30--300 Hz\\(Butterworth 4, causal)};
\node[box, below=of m1, fill=intc!8, draw=intc!70!black] (m2) {TKEO $x_t^2-x_{t-1}x_{t+1}$};
\node[box, below=of m2, fill=intc!8, draw=intc!70!black] (m3) {$|\cdot|$, clip at $5\times$ 99.5th pct., $\sqrt{\cdot}$};
\node[box, below=of m3, fill=intc!8, draw=intc!70!black] (m4) {low-pass 8 Hz\\(Butterworth 4, causal)};
\node[box, below=of m4, fill=intc!8, draw=intc!70!black] (m5) {downsample $\times 8$ to 500 Hz\\(causal; total delay $\approx$ 60 ms)};
% Kinematics column
\node[io, right=6mm of m0, text width=3.9cm] (k0) {kinematics, 36 columns, 500 Hz};
\node[kbox, below=of k0] (k1) {13 features: wrist, index, thumb (9),\\grip aperture, $F_L$, $F_G$, $F_G/F_L$};
\node[kbox, below=of k1] (k2) {hand-near-object gate\\index--object distance 12$\to$20 cm};
\node[kbox, below=of k2] (k3) {$F_L\leftarrow F_L\cdot\mathrm{gate}$,\ \ $F_G\leftarrow F_G\cdot\mathrm{gate}$};
\node[kbox, below=of k3] (k4) {velocity of all 13\\(Savitzky--Golay, 51 samples, order 3)};
\node[kbox, below=of k4] (k5) {26 columns [position 13 | velocity 13]};
\foreach \c in {e,m,k}{\foreach \i/\j in {0/1,1/2,2/3,3/4,4/5}{\draw[arr] (\c\i)--(\c\j);}}
\path let \p1=(e5.south), \p2=(m5.south), \p3=(k5.south) in coordinate (low) at ({\x2},{min(\y1,\y2,\y3)-8mm});
\node[result, anchor=north, text width=12.6cm] (cache) at (low) {cache per series: \texttt{data/cache\_main\_stride500\_square\_pose\_force\_sg51/P\{p\}\_\{split\}\_S\{s\}.npz}\ \ (eeg, kin, emg at 500 Hz, cut to the same length)};
\draw[arr] (e5)--(e5|-cache.north); \draw[arr] (m5)--(cache); \draw[arr] (k5)--(k5|-cache.north);
\end{tikzpicture}
\end{latin}
}
\caption{مرحله‌ی مشترک پیش‌پردازش، از فایل خام تا \lr{cache} ۵۰۰ هرتز. چپ: \lr{EEG}. وسط: \lr{EMG} که پوش آن هدف پیش‌بینی است. راست: سینماتیک و نیروها. همه‌ی گام‌ها روی کل سری انجام می‌شوند، نه روی پنجره‌ها.}
\label{fig:preproc-shared}
\end{figure}

\subsection{مرحله‌ی مشترک: \lr{EEG}}

پنج گام، به همین ترتیب، روی کل سری:
\begin{enumerate}
\item \textbf{\lr{band-pass} \lr{۰٫۵} تا ۴۰ هرتز} (فیلتر \lr{FIR} با پنجره‌ی همینگ، بدون جابه‌جایی زمانی). زیر \lr{۰٫۵} هرتز بیشتر رانش آرام الکترودها و تعریق است و بالای ۴۰ هرتز بیشتر فعالیت عضلات سر و نویز. فعالیت مغزی مرتبط با حرکت (تتا تا گامای پایین) در همین بازه است.
\item \textbf{\lr{notch} ۵۰ هرتز} (\lr{FIR}) برای حذف برق شهر.
\item \textbf{حذف \lr{artifact} با \lr{ICA}.} \lr{ICA} سیگنال ۳۲ کانال را به مؤلفه‌هایی تجزیه می‌کند که از نظر آماری مستقل‌اند. مؤلفه‌هایی که از منبع غیرمغزی می‌آیند حذف و سیگنال دوباره ساخته می‌شود. روش \lr{FastICA} است، برای هر سری جداگانه، با \lr{seed} ثابت ۴۲ و حداکثر ۲۰۰ تکرار. این داده کانال جداگانه‌ای برای حرکت چشم (\lr{EOG}) ندارد. برای همین \lr{Fp1}، که بالای چشم چپ است و پلک زدن را بیشتر از همه ثبت می‌کند، جای کانال چشم را می‌گیرد. هر مؤلفه‌ای که همبستگی‌اش با \lr{Fp1} به‌طور غیرعادی بالا باشد (\lr{z-score} بیشتر از ۲) حذف می‌شود.
\item \textbf{\lr{band-pass} دوم \lr{۰٫۵} تا ۴۰ هرتز.} همان فیلتر گام ۱ یک بار دیگر اجرا می‌شود (در کد نامش «\lr{delta}» است، اما بازه‌اش \lr{۰٫۵} تا ۴۰ هرتز تنظیم شده). این کار هر چیزی را که \lr{ICA} بیرون از این بازه به سیگنال برگردانده پاک می‌کند.
\item \textbf{بریدن مقادیر بسیار بزرگ.} برای هر کانال یک آستانه حساب می‌شود: $\max(5\,\mathrm{MAD}/0.6745,\ 50\,\mu V)$. \lr{MAD} میانه‌ی قدرمطلق فاصله از میانه است و $\mathrm{MAD}/0.6745$ تخمینی از انحراف معیار است که پرت‌ها رویش اثر ندارند. هر مقداری بیرون از $\pm$ آستانه به خود آستانه بریده می‌شود. این کار جهش‌های کوتاه و بزرگ (مثلاً از حرکت سیم) را محدود می‌کند، بدون اینکه بقیه‌ی سیگنال تغییر کند.
\end{enumerate}
مرجع مشترک (\lr{CAR}) اعمال نمی‌شود و سیگنال با همان مرجع ثبت باقی می‌ماند.

\paragraph{چه چیزی ممکن است باقی بماند.} \lr{ICA} با این روش بیشتر پلک و حرکت چشم را حذف می‌کند. فعالیت عضلات سر و گردن که در فرکانس‌های بالاتر است، و حرکت چشمی که با \lr{Fp1} همبستگی کمی دارد، ممکن است باقی بمانند. این همان هشداری است که در فصل‌های تفسیر درباره‌ی سهم الکترودهای پیشانی و پس‌سری و باند گامای پایین آمده است.

\subsection{مرحله‌ی مشترک: \lr{EMG} و ساختن پوش}

هدف هر دو مدل پوش \lr{EMG} است: منحنی آرامی که نشان می‌دهد عضله در هر لحظه چقدر فعال است. پنج گام:
\begin{enumerate}
\item \textbf{\lr{band-pass} ۳۰ تا ۳۰۰ هرتز} (\lr{Butterworth} مرتبه‌ی ۴). زیر ۳۰ هرتز بیشتر \lr{artifact} حرکت کابل و پوست است و بالای ۳۰۰ هرتز نویز.
\item \textbf{\lr{TKEO}:} $\psi_t=x_t^2-x_{t-1}x_{t+1}$. این عملگر «انرژی» لحظه‌ای سیگنال را می‌دهد. برای نوسان‌های سریع و بزرگ (فعالیت عضله) بزرگ است و برای نویز آرام کوچک؛ برای همین شروع فعالیت را تیزتر از یکسوسازی ساده نشان می‌دهد.
\item \textbf{قدرمطلق، بریدن و جذر.} قدرمطلق گرفته می‌شود و هر مقدار بیشتر از ۵ برابر صدک \lr{۹۹٫۵} همان کانال به این حد بریده می‌شود، تا چند جهش نادر بقیه‌ی سیگنال را کوچک نکنند. سپس جذر گرفته می‌شود، چون \lr{TKEO} واحد توان دارد و جذر آن را به واحد دامنه برمی‌گرداند.
\item \textbf{\lr{low-pass} ۸ هرتز} (\lr{Butterworth} مرتبه‌ی ۴). نوسان‌های سریع حذف می‌شوند و فقط پوش آرام می‌ماند.
\item \textbf{کاهش نرخ از ۴۰۰۰ به ۵۰۰ هرتز} (ضریب ۸، با فیلتر ضدتداخل)، تا با \lr{EEG} و سینماتیک هم‌نرخ شود.
\end{enumerate}

\paragraph{تأخیر پوش.} همه‌ی فیلترهای \lr{EMG} علّی‌اند؛ یعنی فقط از گذشته استفاده می‌کنند و هیچ‌وقت از آینده‌ی سیگنال خبر ندارند. بهای این کار تأخیر است: پوش حدود ۶۰ میلی‌ثانیه دیرتر از فعالیت واقعی عضله بالا و پایین می‌رود. بیشتر این تأخیر از \lr{low-pass} است (۵۲ تا ۶۵ میلی‌ثانیه، بسته به فرکانس) و حدود ۵ میلی‌ثانیه از کاهش نرخ. در مقابل، فیلترهای \lr{EEG} و سینماتیک بدون تأخیرند. پس هدف کمی عقب‌تر از ورودی‌هاست. مدل پایه این را با تأخیر آموختنی خودش جبران می‌کند، و \lr{v2} با دیدن ۸ ثانیه بافت می‌تواند ورودی‌های کمی قبل‌تر را هم به کار ببرد.

\subsection{مرحله‌ی مشترک: سینماتیک و نیروها}

\paragraph{سیزده ویژگی.} از ۳۶ ستون خام، ۱۳ ویژگی ساخته می‌شود:
\begin{itemize}
\item مکان سه‌بعدی مچ، نوک انگشت اشاره و نوک شست (۹ ویژگی، سانتی‌متر)؛
\item دهانه‌ی گرفتن: فاصله‌ی نوک اشاره تا نوک شست (۱ ویژگی)؛
\item نیروی بار $F_L$: اندازه‌ی بردار نیروی سه‌بعدی صفحه‌ی اول، $\sqrt{F_x^2+F_y^2+F_z^2}$؛
\item نیروی گرفتن $F_G$: میانگین قدرمطلق نیروی عمودی بر دو صفحه، $(|F_{z1}|+|F_{z2}|)/2$؛
\item نسبت $F_G/(F_L+1)$.
\end{itemize}
ده ویژگی اول «وضعیت دست» را توصیف می‌کنند و سه ویژگی آخر نیروها را.

\paragraph{\lr{gate} نزدیکی دست به جسم.} بین دو بلندکردن، آزمایشگر جسم را جابه‌جا یا سطح آن را عوض می‌کند و صفحه‌های نیروسنج این را هم ثبت می‌کنند. در \lr{P1}، در سری‌هایی که سطح عوض می‌شود، صفحه‌ها در حدود ۱۹ درصد زمان سری بار ۲ تا ۵ نیوتنی دارند، در حالی که دست شرکت‌کننده حدود ۲۸ سانتی‌متر دورتر است و عضله‌ها در استراحت‌اند. اگر این نیرو به مدل داده شود، مدل در استراحت جهش \lr{EMG} پیش‌بینی می‌کند.

راه حل یک \lr{gate} است که از فاصله‌ی نوک انگشت اشاره تا جسم ساخته می‌شود. این فاصله دو حالت روشن دارد: حدود \lr{۱۰٫۵} سانتی‌متر وقتی جسم در دست است و حدود ۲۸ سانتی‌متر وقتی دست کنار است. \lr{gate} تا فاصله‌ی ۱۲ سانتی‌متر برابر ۱ است، تا ۲۰ سانتی‌متر به‌طور خطی به صفر می‌رسد و بعد از آن صفر است. اگر حسگر مکان را گم کرده باشد، \lr{gate} صفر در نظر گرفته می‌شود. سپس \lr{gate} با پنجره‌ی \lr{۰٫۱} ثانیه نرم می‌شود: اول بیشینه‌ی متحرک، بعد میانگین متحرک، هر دو متقارن. ترتیب این دو مهم است: بیشینه‌ی متحرک تضمین می‌کند که نرم کردن هیچ‌وقت \lr{gate} را در ناحیه‌ی نزدیک کمتر از ۱ نکند، حتی وقتی دست سریع به جسم می‌رسد. در آخر $F_L$ و $F_G$ در \lr{gate} ضرب می‌شوند.

\paragraph{سرعت.} سرعت هر ۱۳ ویژگی با مشتق اول فیلتر \lr{Savitzky-Golay} حساب می‌شود (پنجره‌ی ۵۱ نمونه یعنی ۱۰۲ میلی‌ثانیه، چندجمله‌ای مرتبه‌ی ۳). این فیلتر در هر پنجره یک چندجمله‌ای برازش می‌کند و مشتق آن را می‌گیرد، که بسیار کم‌نویزتر از تفاضل ساده‌ی دو نمونه است. پنجره‌ی ۵۱ نمونه‌ای لرزش حسگر را در زمان استراحت کم می‌کند. سرعت نیروها بعد از اعمال \lr{gate} حساب می‌شود، پس وقتی دست دور است سرعت نیرو هم صفر است. شتاب حساب نمی‌شود، چون مشتق دوم لرزش حسگر را بسیار بزرگ می‌کند.

نتیجه ۲۶ ستون است: ۱۳ ویژگی و ۱۳ سرعت. هر دو مدل نسبت نیروها و سرعتش را کنار می‌گذارند و ۲۴ ویژگی باقی می‌ماند: وضعیت دست (۱۰)، دو نیروی \lr{gated}، سرعت وضعیت دست (۱۰) و سرعت دو نیرو.

\paragraph{\lr{cache}.} برای هر سری، سه سیگنال ۵۰۰ هرتز به طول کوتاه‌ترینشان بریده و در پوشه‌ی \lr{cache} ذخیره می‌شوند (\lr{data/cache\_main\_\allowbreak stride500\_\allowbreak square\_\allowbreak pose\_\allowbreak force\_\allowbreak sg51}). \lr{v2} فقط این فایل‌ها را می‌خواند و خودش هیچ‌کدام از گام‌های بالا را دوباره اجرا نمی‌کند. پس دو مدل دقیقاً از یک داده‌ی تمیزشده شروع می‌کنند و تفاوت نتایج از مدل‌هاست، نه از پیش‌پردازش.

\subsection{مرحله‌ی ویژه‌ی \lr{v2}: از \lr{cache} تا \lr{tokens}}
\label{sec:preproc-v2}

شکل~\ref{fig:preproc-v2} این مرحله را نشان می‌دهد. یک اصل در همه‌ی گام‌ها رعایت می‌شود: هر فیلتری اول روی کل سری با ۵۰۰ هرتز اجرا می‌شود و بعد نرخ کم می‌شود. فیلترها برای کار درست به سیگنال طولانی نیاز دارند؛ اگر روی پنجره‌های کوتاه اجرا شوند، لبه‌های هر پنجره خراب می‌شود.

\begin{figure}[p]
\centering
\resizebox{\textwidth}{!}{% v2 preprocessing: 500 Hz cache -> 50 Hz model inputs and target.
\begin{latin}
\begin{tikzpicture}[node distance=4mm, font=\footnotesize, box/.append style={text width=2.85cm}]
\node[io, text width=13cm] (c) {cache of one series at 500 Hz (eeg $T\times32$, kin $T\times26$, emg $T\times5$)};
% kinematics
\node[kbox, anchor=north] (k1) at ($(c.south)+(-5.04,-0.8)$) {drop $F_G/F_L$ and its velocity: 26 $\to$ 24};
\node[kbox, below=of k1] (k2) {to tokens: FIR anti-alias, zero-phase, $\div 10$};
\node[kbox, below=of k2] (k3) {z-score with train series};
\node[result, below=of k3, text width=2.85cm] (k4) {$X^{k}$\ \ $T_{50}\times24$};
% history
\node[kbox, right=3mm of k1] (h1) {gated $F_L$, $F_G$ (500 Hz)};
\node[kbox, below=of h1] (h2) {holds: $F_L>1$ N for $\ge0.5$ s};
\node[kbox, below=of h2] (h3) {8 features, each known only after the hold ends};
\node[kbox, below=of h3] (h3b) {every 10th sample (no filter)};
\node[result, below=of h3b, text width=2.85cm] (h4) {$X^{h}$\ \ $T_{50}\times8$};
% EEG
\node[ebox, right=3mm of h1] (e1) {4 bands: $\theta,\alpha,\beta,\gamma_{low}$\\Butterworth 4, zero-phase};
\node[ebox, below=of e1] (e2) {power $|\mathrm{Hilbert}|^2$, low-pass 5 Hz};
\node[ebox, below=of e2] (e3) {$\log\max(P,\ \epsilon)$\\$\epsilon=10^{-3}\,\mathrm{median}\,P$};
\node[ebox, below=of e3] (e3b) {to tokens ($\div 10$), z-score, clip $\pm5$};
\node[result, below=of e3b, text width=2.85cm] (e4) {$X^{e}$\ \ $T_{50}\times32\times4$};
% EMG
\node[box, right=3mm of e1, fill=intc!8, draw=intc!70!black] (m1) {EMG envelope (500 Hz)};
\node[box, below=of m1, fill=intc!8, draw=intc!70!black] (m2) {to tokens ($\div 10$), $\max(\cdot,0)$};
\node[box, below=of m2, fill=intc!8, draw=intc!70!black] (m3) {$s_c$ = 95th pct. of train tokens};
\node[box, below=of m3, fill=intc!8, draw=intc!70!black] (m3b) {$u=\sqrt{y/s_c}$};
\node[result, below=of m3b, text width=2.85cm] (m4) {target $u$\ \ $T_{50}\times5$};
\foreach \c in {k,h,e,m}{\draw[arr] (c.south-|\c1)--(\c1);}
\foreach \a/\b in {k1/k2,k2/k3,k3/k4,h1/h2,h2/h3,h3/h3b,h3b/h4,e1/e2,e2/e3,e3/e3b,e3b/e4,m1/m2,m2/m3,m3/m3b,m3b/m4}{\draw[arr] (\a)--(\b);}
\node[result, below=7mm of $(h4.south)!0.5!(e4.south)$, text width=13cm] (w) {training windows: 400 tokens (8 s), stride 100 tokens (2 s), random start offset in $[0,100)$ each epoch\\test: windows of 8 s every 4 s, keep the middle 4 s, spline back to 500 Hz};
\foreach \x in {k4,h4,e4,m4}{\draw[arr] (\x)--(\x|-w.north);}
\end{tikzpicture}
\end{latin}
}
\caption{مرحله‌ی ویژه‌ی \lr{v2}، از \lr{cache} ۵۰۰ هرتز تا ورودی‌ها و هدف ۵۰ هرتز. از چپ: سینماتیک، تاریخچه‌ی بلندکردن‌ها (فقط از نیروهای \lr{gated})، \lr{band power} \lr{EEG} و هدف \lr{EMG}. کادر پایین: ساختن پنجره‌های آموزش و پیش‌بینی روی سری آزمون.}
\label{fig:preproc-v2}
\end{figure}

\subsubsection*{۱. تبدیل به \lr{tokens} (کاهش نرخ از ۵۰۰ به ۵۰ هرتز)}
اگر فقط هر دهمین نمونه را برداریم، نوسان‌های بالای ۲۵ هرتز (نصف نرخ جدید) از بین نمی‌روند و به شکل نوسان‌های آرام جعلی در سیگنال جدید ظاهر می‌شوند (\lr{aliasing}). برای همین اول یک فیلتر ضدتداخل اجرا می‌شود: فیلتر \lr{FIR} با ۲۰۱ ضریب (\lr{scipy.signal.decimate})، رفت و برگشتی تا سیگنال در زمان جابه‌جا نشود. \lr{token} شماره‌ی $k$ دقیقاً روی نمونه‌ی $10k$ قرار می‌گیرد.

یک مشکل در لبه‌ها هست: این فیلتر بیرون سری را صفر فرض می‌کند، پس اولین و آخرین \lr{tokens} تا حدود نصف مقدار واقعی‌شان به سمت صفر کشیده می‌شوند. برای رفع آن، سری پیش از فیلتر در هر دو سر ۱۰۰ نمونه با «بازتاب فرد» ادامه داده می‌شود: $x_{-i}=2x_0-x_i$. با این بازتاب، سیگنال ثابت ثابت می‌ماند و سیگنالی که خطی بالا می‌رود همچنان خطی ادامه پیدا می‌کند. پس از فیلتر، بخش اضافه دوباره بریده می‌شود. سری باید دست‌کم ۷۰۰ نمونه داشته باشد تا این فیلتر کار کند.

\subsubsection*{۲. سینماتیک}
۲۴ ستون سینماتیک (بعد از کنار گذاشتن نسبت نیروها) به \lr{tokens} تبدیل می‌شوند. سپس هر ویژگی با میانگین و انحراف معیار \lr{tokens} سری‌های آموزش استاندارد می‌شود (\lr{z-score}). این ویژگی‌ها بریده نمی‌شوند، چون \lr{gate} و فیلترهای مرحله‌ی مشترک مقادیر پرت آن‌ها را از قبل گرفته‌اند. نتیجه $X^{k}\in\mathbb{R}^{T\times24}$ است.

\subsubsection*{۳. تاریخچه‌ی بلندکردن‌ها}
این ویژگی‌ها فقط از دو نیروی \lr{gated} و با نرخ ۵۰۰ هرتز ساخته می‌شوند؛ \lr{EMG} هیچ نقشی در آن‌ها ندارد.

\textbf{پیدا کردن نگه‌داشتن‌ها.} یک «نگه‌داشتن» بازه‌ای پیوسته است که در آن $F_L$ بیشتر از ۱ نیوتن است و دست‌کم \lr{۰٫۵} ثانیه طول می‌کشد. چون $F_L$ \lr{gated} است، نیرویی که آزمایشگر وارد می‌کند نگه‌داشتن حساب نمی‌شود.

\textbf{بدون نگاه به آینده.} مدل فقط وقتی از یک نگه‌داشتن خبردار می‌شود که آن نگه‌داشتن تمام شده باشد. همه‌ی ویژگی‌های مربوط به نگه‌داشتن $k$ از نمونه‌ی پایان آن به بعد مقدار می‌گیرند. پس مقدار هر ویژگی در زمان $t$ فقط به نمونه‌های تا $t$ بستگی دارد.

\textbf{هشت ویژگی.} ویژگی‌ها با عددهای ثابت مقیاس می‌شوند، نه با \lr{z-score}، تا معنای هر عدد روشن بماند و مقدارها تقریباً بین ۰ و ۱ باشند:
\begin{center}\small
\begin{tabular}{rll}
\toprule
& ویژگی & مقیاس \\
\midrule
۱ & آیا تا اینجا نگه‌داشتنی بوده است & ۰ یا ۱ \\
۲ & تعداد نگه‌داشتن‌ها تا اینجا & تقسیم بر ۴۰ \\
۳ & زمان از شروع سری & تقسیم بر ۳۰۰ ثانیه \\
۴ & زمان از پایان آخرین نگه‌داشتن & حداکثر ۳۰ ثانیه، تقسیم بر ۳۰ \\
۵ و ۶ & میانگین $F_L$ و $F_G$ در آخرین نگه‌داشتن & تقسیم بر ۱۰ نیوتن \\
۷ و ۸ & میانگین $F_L$ و $F_G$ در همه‌ی نگه‌داشتن‌های قبلی & تقسیم بر ۱۰ نیوتن \\
\bottomrule
\end{tabular}
\end{center}
ویژگی ۳ از ابتدای سری مقدار دارد و بقیه تا پایان اولین نگه‌داشتن صفرند.

\textbf{تبدیل به \lr{tokens} بدون فیلتر.} این ویژگی‌ها پله‌ای‌اند: تا پایان یک نگه‌داشتن ثابت می‌مانند و ناگهان عوض می‌شوند. برای همین از هر ۱۰ نمونه فقط یکی برداشته می‌شود و فیلتر ضدتداخل اجرا نمی‌شود. فیلتر رفت و برگشتی، پله را در هر دو طرف پخش می‌کند و تغییر را چند \lr{token} زودتر از پایان واقعی نگه‌داشتن نشان می‌دهد؛ یعنی کمی از آینده به مدل خبر می‌دهد. نتیجه $X^{h}\in\mathbb{R}^{T\times8}$ است.

همین شماره‌گذاری نگه‌داشتن‌ها برای هر نمونه هم نگه داشته و به همین روش به \lr{tokens} تبدیل می‌شود. جزء \lr{lift} در \lr{loss function} از آن استفاده می‌کند تا میانگین پیش‌بینی و هدف را در هر نگه‌داشتن مقایسه کند.

\subsubsection*{۴. \lr{band power} \lr{EEG}}
\lr{v2} به‌جای موج خام \lr{EEG}، توان آن را در چهار باند می‌بیند. موج خام \lr{۰٫۵} تا ۴۰ هرتز تقریباً هیچ اطلاعات خطی درباره‌ی پوش \lr{EMG} ندارد، اما تغییرهای آرام توان باندها دارد؛ مثلاً کاهش توان بتا روی قشر حرکتی پیش از حرکت (\lr{ERD}). برای هر باند تتا (۴ تا ۸)، آلفا (۸ تا ۱۳)، بتا (۱۳ تا ۳۰) و گامای پایین (۳۰ تا ۴۰ هرتز)، و برای هر الکترود، روی کل سری:
\begin{enumerate}
\item \textbf{جدا کردن باند:} \lr{band-pass} \lr{Butterworth} مرتبه‌ی ۴، رفت و برگشتی.
\item \textbf{توان لحظه‌ای:} مجذور دامنه‌ی سیگنال تحلیلی، $P_t=|\mathcal{H}(x)_t|^2$ (تبدیل هیلبرت). این عدد اندازه‌ی نوسان را در هر لحظه می‌دهد، مستقل از اینکه موج در آن لحظه در قله باشد یا در صفر.
\item \textbf{نرم کردن:} \lr{low-pass} ۵ هرتز. فیلتر مرتبه‌ی ۱ است و رفت و برگشتی اجرا می‌شود، طوری طراحی شده که کل فیلتر دقیقاً در ۵ هرتز \lr{−۳} دسی‌بل کاهش بدهد. مرتبه‌ی ۱ عمداً انتخاب شده: پاسخ ضربه‌ی آن همیشه مثبت است، پس توان نرم‌شده هیچ‌وقت منفی نمی‌شود. با مرتبه‌ی ۴، توان نرم‌شده در \lr{۰٫۳} تا \lr{۲٫۴} درصد نمونه‌ها منفی می‌شد و بعد از \lr{log} جهش‌های بزرگ می‌ساخت.
\item \textbf{لگاریتم با کف:} $\log\max(P,\ 10^{-3}\cdot\operatorname{median}P)$، که میانه برای هر الکترود و باند جداگانه روی کل سری حساب می‌شود. توان توزیع بسیار کجی دارد و لگاریتم آن را متقارن‌تر می‌کند؛ علاوه بر این، تغییر نسبی توان (مثلاً «۲۰ درصد کاهش») در لگاریتم به یک اختلاف ثابت تبدیل می‌شود. کف نمی‌گذارد یک لحظه‌ی تقریباً بی‌سیگنال، بعد از لگاریتم به عدد منفی بسیار بزرگ تبدیل شود.
\end{enumerate}
نتیجه برای هر لحظه یک جدول $32\times4$ است. این جدول به \lr{tokens} تبدیل می‌شود، هر یک از ۱۲۸ خانه‌ی آن با آماره‌های \lr{tokens} آموزش استاندارد می‌شود و مقادیر بیرون از $\pm5$ انحراف معیار بریده می‌شوند. بریدن جلوی اثر چند لحظه‌ی \lr{artifact} باقی‌مانده را می‌گیرد. نتیجه $X^{e}\in\mathbb{R}^{T\times32\times4}$ است و الکترودها به همان ترتیب \lr{hs.eeg.names} می‌مانند تا \lr{electrode graph} بتواند فاصله‌ی آن‌ها روی سر را بداند.

\subsubsection*{۵. هدف}
پوش ۵۰۰ هرتز \lr{EMG} با همان روش به \lr{tokens} تبدیل می‌شود. فیلتر ضدتداخل گاهی کنار جهش‌های تیز مقدار کمی زیر صفر می‌سازد؛ این مقادیر صفر می‌شوند، چون شدت فعالیت نمی‌تواند منفی باشد. برای هر عضله $s_c$ برابر صدک ۹۵ \lr{tokens} آموزش است، و هدف مدل $u=\sqrt{y/s_c}$ است (بخش~\ref{sec:v2}). تقسیم بر $s_c$ پنج عضله را هم‌مقیاس می‌کند، چون دامنه‌ی خام \lr{EMG} از عضله‌ای به عضله‌ی دیگر بسیار فرق دارد.

ارزیابی روی این هدف ۵۰ هرتز انجام نمی‌شود. پیش‌بینی با درون‌یابی اسپلاین به ۵۰۰ هرتز برمی‌گردد، به واحد اصلی تبدیل می‌شود ($\hat y=s_c\hat u^2$) و با پوش اصلی ۵۰۰ هرتز مقایسه می‌شود. پس هر خطایی که کاهش نرخ ایجاد کند در اعداد گزارش‌شده هم دیده می‌شود.

\subsubsection*{۶. پنجره‌ها}
ورودی‌ها برای هر سری یک بار و کامل نرمال می‌شوند و بعد پنجره‌ها از آن بریده می‌شوند. پس همه‌ی پنجره‌های یک سری با یک مقیاس نرمال شده‌اند.
\begin{itemize}
\item \textbf{آموزش:} پنجره‌های ۴۰۰ \lr{token} (۸ ثانیه) با فاصله‌ی ۱۰۰ \lr{token} (۲ ثانیه)، که هر دو پنجره‌ی پشت سر هم ۶ ثانیه هم‌پوشانی دارند. در هر دوره، نقطه‌ی شروع اولین پنجره به‌طور تصادفی بین ۰ و ۹۹ \lr{token} انتخاب می‌شود. بنابراین مدل در هر دوره برش‌های کمی متفاوتی از داده می‌بیند و لبه‌ی پنجره‌ها همیشه روی یک لحظه‌ی ثابت از سری نمی‌افتد.
\item \textbf{\lr{validation} و آزمون:} پنجره‌های ۸ ثانیه‌ای هر ۴ ثانیه، با نگه داشتن ۴ ثانیه‌ی وسط هر پنجره (شکل~\ref{fig:stitch}).
\end{itemize}

\subsection{مقایسه‌ی دو پیش‌پردازش}

\begin{table}[htbp]
\centering\footnotesize
\caption{تفاوت پیش‌پردازش مدل پایه و \lr{v2}. مرحله‌ی مشترک (تا \lr{cache}) در هر دو یکسان است.}
\label{tab:preproc}
\begin{tabular}{p{0.2\textwidth}p{0.33\textwidth}p{0.37\textwidth}}
\toprule
& مدل پایه & \lr{v2} \\
\midrule
نرخ ورودی مدل & ۵۰۰ هرتز & ۵۰ هرتز (\lr{tokens}، با فیلتر ضدتداخل) \\
طول پنجره & ۱ ثانیه، بدون هم‌پوشانی & ۸ ثانیه، با فاصله‌ی ۲ ثانیه و شروع تصادفی \\
\lr{EEG} & موج زمانی \lr{۰٫۵} تا ۴۰ هرتز، ۳۲ کانال، بدون نرمال‌سازی بیشتر & لگاریتم \lr{band power} در ۴ باند، $32\times4$ \\
سینماتیک & ۲۴ ویژگی، \lr{z-score} & همان ۲۴ ویژگی، \lr{z-score} پس از کاهش نرخ \\
تاریخچه & ندارد & ۸ ویژگی از نیروهای \lr{gated}، بدون نگاه به آینده \\
هدف & پوش، \lr{z-score} با میانگین و انحراف معیار آموزش & $\sqrt{y/s_c}$ با $s_c$ صدک ۹۵ آموزش \\
بریدن & -- & \lr{EEG} در $\pm5$؛ هدف در صفر از پایین \\
ارزیابی & روی پنجره‌های ۱ ثانیه‌ای & روی کل سری، پس از برگرداندن به ۵۰۰ هرتز \\
\bottomrule
\end{tabular}
\end{table}

جدول~\ref{tab:preproc} تفاوت‌ها را خلاصه می‌کند. همه‌ی تصمیم‌های \lr{v2} در این مرحله یکی از سه هدف را دنبال می‌کنند: دیدن بافت طولانی‌تر با هزینه‌ی کم (\lr{tokens} ۵۰ هرتز)، دادن اطلاعاتی به مدل که سینماتیک لحظه‌ای ندارد (تاریخچه و \lr{band power})، و جلوگیری از نشت اطلاعات (آماره‌ها فقط از آموزش، تاریخچه فقط از گذشته، تاریخچه بدون فیلتر).
